\chapter{Supreme Court on Child Welfare, Fundamental Rights and Bail}
\label{ch:sc-child-welfare-fundamental-rights-bail}

This lecture surveys the Supreme Court judgments delivered over roughly the past one
year, organised around a single unifying theme: the \emph{constitutionalization} of
important subjects. Just as the Court has earlier taken an expansive view of women's
rights and the environment, it now reads child rights into the Constitution — and, more
broadly, keeps widening the scope of \textbf{Article 21} and the fundamental rights.
The lecture is divided into four substantive blocks: child welfare, recent judgments
upholding fundamental rights, online takedowns and blocking, and the jurisprudence of
bail.

\section{Supreme Court and Child Welfare}
\label{sec:child-welfare}

Child welfare is approached in three layers: first the \emph{constitutional}
provisions, then the \emph{statutory} (legal) framework, and finally the Supreme Court's
judgments — split into those that are \emph{historically} important and those that are
\emph{recent}.

\subsection{Constitutional provisions}

The constitutional jurisprudence on children can be divided into Fundamental Rights,
Directive Principles of State Policy (DPSPs), and one Fundamental Duty.

\begin{KeyConcept}
The theme of the entire lecture is \textbf{constitutionalization}. When the Court speaks
of an \emph{expansive} or \emph{broad} view under Article 21, it means that the
Fundamental Rights, DPSPs and Fundamental Duties already in the Constitution are
treated as the base framework which the Court then \emph{expands} to cover a new subject
(here, child rights) — rather than the Constitution being read narrowly.
\end{KeyConcept}

\textbf{Fundamental Rights.} The relevant Fundamental Rights are:

\begin{itemize}
  \item \textbf{Article 14} — the general guarantee of equality before law.
  \item \textbf{Article 15} — permits special provisions for \emph{women and children}.
  \item \textbf{Article 21} — life and personal liberty; it applies to everyone,
        and therefore to children as well.
  \item \textbf{Article 21A} — free and compulsory education for children aged 6 to 14.
  \item \textbf{Articles 23 and 24} — together called the \emph{right against
        exploitation}.
\end{itemize}

\begin{CriticalPoint}
Most of the Supreme Court judgments that have spoken about the rights of the child
revolve around \textbf{Articles 23 and 24} (the right against exploitation), which is
why this pair carries outsized weight in child-rights jurisprudence.
\end{CriticalPoint}

\textbf{Directive Principles of State Policy.} The DPSPs relevant to children are:

\begin{itemize}
  \item \textbf{Article 39(e)} — protection of children from abuse.
  \item \textbf{Article 39A} — free legal aid; important for child-related cases, since
        children need to be given the benefit of free legal aid and free lawyers for the
        family.
  \item \textbf{Article 45} — early childhood care and education for children below six
        years.
  \item \textbf{Article 47} — public health and nutrition.
\end{itemize}

\textbf{Fundamental Duty.} The relevant Fundamental Duty is the duty of a \emph{parent
or guardian} to provide opportunities for education to a child between the ages of 6 and
14 years.

\subsection{Statutory and legal framework}

Alongside the Constitution, there is a \emph{statutory} (legal) framework of laws that
deal with the rights of the child. These are the statutes an examiner expects you to
cite alongside the constitutional provisions.

\begin{GovtPolicy}
The key statutes in the child-welfare framework are:
\begin{itemize}
  \item \textbf{Immoral Traffic (Prevention) Act, 1956} (ITPA).
  \item \textbf{Juvenile Justice Act, 2005} — children who commit crimes are kept before
        the Juvenile Justice Board rather than tried as adults; this is where the concern
        over the \emph{delinquency rate} is discussed.
  \item \textbf{Criminal Law (Amendment) Act, 2013} — enacted in the aftermath of the
        Nirbhaya case, where a juvenile was involved; it also deals with trafficking.
  \item \textbf{Bharatiya Nyaya Sahita} — carries provisions on the \emph{selling and
        buying of minors}.
  \item \textbf{Protection of Children from Sexual Offences (POCSO) Act, 2012} —
        \emph{Prevention of Children from Sexual Offences}; it set up \textbf{Fast Track
        Courts}.
\end{itemize}
\end{GovtPolicy}

\textbf{The MTP--POCSO overlap.} The lecture flags a sharp, contested overlap between
the \textbf{Medical Termination of Pregnancy (MTP) Act, 1971} and POCSO. At one end you
want to promote the reproductive rights of a woman carrying a pregnancy; at the other,
POCSO protects the child.

\begin{CriticalPoint}
The underlying tension is between the \emph{rights of the fetus} and the \emph{rights of
the woman}. The Supreme Court has itself given controversial opinions: at one point
Justice Chandrachud spoke of the right of the fetus as decisive; at another he held that
the \emph{reproductive choice / autonomy} of the woman is what is guaranteed — not only
under Article 21 but explicitly recognised in \textbf{K.S. Puttaswamy} (right to privacy
as a fundamental right). The balance between these two is what makes the area
controversial.
\end{CriticalPoint}

\begin{ExampleBox}
The MTP Act permits termination only up to a gestational limit, though for certain
specific categories of women the period is extended \emph{from 20 to 24 weeks}. A rape
victim, for instance, is both a victim of an offence triable under POCSO \emph{and} a
person with higher grounds for termination under the MTP Act. The practical trap arises
in cases of adolescent (e.g.\ 18-year-old and 17-year-old) consensual relations: because
a doctor is bound to treat abortion as an exception to criminal law and must inform the
police, the male adolescent can end up within the POCSO framework even where the act was
consensual and no complaint is filed. India was one of the most progressive nations to
enact an anti-abortion-regulating law as early as 1971, but between 1971 and 2026 the
jurisprudence of abortion and bodily privacy has changed so much that the statutory
framework is widely seen as needing amendment.
\end{ExampleBox}

\subsection{Historically important judgments}

The lecture traces the history of child-rights cases through a series of landmark
judgments, several of which began in the PIL (public interest litigation) era.

\begin{itemize}
  \item \textbf{Bandhua Mukti Morcha (1984)} — a landmark PIL-era case that expanded
        Article 21 into a guarantee of \emph{living in dignified conditions} for bonded
        labourers and vulnerable workers; because child labour is implicated, it carries
        an automatic ramification for child rights.
  \item \textbf{Lakshmi Kant Pandey (1984)} — the foundational case on \emph{inter-country
        adoption}; the Court laid down safeguards to prevent the \emph{commercialization
        of adoption} and the \emph{exploitation of abandoned children}.
  \item \textbf{Unni Krishnan (1993)} — held that the right to education is implicit in
        Article 21, \emph{before} Article 21A was inserted; the Court read Article 21
        \emph{with} the DPSPs — a classic example of the \emph{harmonious interpretation}
        of Fundamental Rights and DPSPs.
  \item \textbf{M.C. Mehta v. State of Tamil Nadu (1996)} — laid down guidelines
        prohibiting the \emph{employment of children in hazardous industries}.
  \item \textbf{Gaurav Jain (1997)} — concerning children of sex workers; the Court
        rejected segregated schools and hostels for such children and recommended
        \emph{mainstreaming}, care, protection and rehabilitation.
  \item \textbf{Sakshi (2004)} — issued child-friendly procedural directions for recording
        the evidence of child victims of sexual abuse, so that children are not exposed to
        a vulnerable environment while giving evidence.
  \item \textbf{Bachpan Bachao Andolan} — directions to address the widespread
        exploitation and trafficking of children.
  \item \textbf{Sampurna Behrua (2018)} — wide-ranging directions for the Juvenile Justice
        Board, the Child Welfare Committee and Child Care Institutions.
\end{itemize}

\begin{GovtPolicy}
\textbf{CARA — Central Adoption Resource Authority.} The Lakshmi Kant Pandey judgment was
the foundation for \textbf{CARA}, the organisation that manages the adoption of a child
and functions under the \textbf{Ministry of Women and Child Development}. It is an
example of \emph{institutional convergence} — coordination between Parliament and the
judiciary — and can be cited under the doctrine of separation of powers: the Court's PIL
was followed by statutory creation of the authority. The same case was also a foundation
for the juvenile-justice framework.
\end{GovtPolicy}

\subsection{Recent judgments}

The recent (past-year) child-welfare decisions include:

\begin{itemize}
  \item \textbf{Nipun Saxena (2000, as stated in lecture)} — protection of the identity of
        rape survivors and child sexual-abuse victims.
  \item \textbf{Right to menstrual health (2026)} — the Supreme Court ruled that the right
        to menstrual health is a fundamental right guaranteed under \textbf{Article 21}.
  \item \textbf{Hamsa Nandhini} — the Court upheld the rights of \emph{adoptive mothers}
        under the \textbf{Maternity Benefit Act, 1961}, and thereby the rights of the
        adopted child.
  \item \textbf{K.P. Kiran Kumar (2026)} — strict guidelines to prevent \emph{child
        trafficking}.
\end{itemize}

\begin{PYQLink}
A question the Court can pose is: \emph{discuss the expanding jurisprudence of social
justice in recent times, with examples from recent case law}. The child-welfare cases
above are the ready-made illustration for that question.
\end{PYQLink}

\begin{QuickRevision}
Child welfare is read through \emph{constitutional} provisions (Articles 14, 15, 21,
21A, 23--24; DPSPs 39(e), 39A, 45, 47; the parent's Fundamental Duty), a \emph{statutory}
framework (ITPA 1956, JJ Act 2005, Criminal Law Amendment Act 2013, BNS, POCSO 2012, and
the MTP--POCSO overlap), and the Court's own judgments — historically Bandhua Mukti Morcha
(1984), Lakshmi Kant Pandey (1984), Unni Krishnan (1993), M.C. Mehta (1996), Gaurav Jain
(1997), Sakshi (2004), Bachpan Bachao Andolan, Sampurna Behrua (2018); recently Nipun
Saxena, menstrual health (2026), Hamsa Nandhini and K.P. Kiran Kumar (2026).
\end{QuickRevision}

\section{Recent Supreme Court Decisions Upholding Fundamental Rights}
\label{sec:recent-fundamental-rights}

A recurring pattern of the past year is the Court's continuing \emph{expansion of
Article 21} into new trajectories. Most of these cases are grounded in Article 21 (a few
also invoke Article 19), and they are the examples to cite when a question asks how the
framework of fundamental rights has been expanded by the Court.

\begin{itemize}
  \item \textbf{Amar Jain} — the Court held Article 21 to encompass the \emph{right to
        digital access}, and directed companies to revise their \textbf{KYC norms} so
        that physically disabled persons are not denied access to banking facilities
        through facial-recognition and similar requirements.
  \item \textbf{Rajeev Raturi (2024)} — the Court ruled that \emph{accessibility} is
        central to the \emph{right to life} under Article 21.
  \item \textbf{Save Life Foundation (2026)} — \emph{right to trauma care} is integral to
        Article 21; it extends from the site of injury to definitive hospital treatment
        and covers the entire chain that connects an injured person to medical care
        (paramedic, nurse, and the receiving facility).
  \item \textbf{Peela Pahan (as stated in lecture)} — the Court ruled that Article 21 is
        not confined to \emph{speedy trial} alone; the \emph{timely pronouncement of
        reserved judgments} is itself a part of Article 21.
  \item \textbf{Maniyar Iliyasi (as stated in lecture)} — the \emph{right to walk on a
        demarcated footpath} is a fundamental right guaranteed under Articles 19 and 21.
  \item \textbf{Right to safe travel on highways} — the right to safe travel on highways
        is a part of the right to life; the Court issued detailed directions.
  \item \textbf{Right to move freely in public places without fear of attack} — a
        two-judge bench (including Justice Pardiwala) interpreted the \emph{ABC rules}
        under the \textbf{Prevention of Cruelty to Animals Act, 1960} to hold that stray
        dogs must be picked up, castrated, and kept restricted in animal centres, so that
        they do not curtail a person's fundamental right to move freely without fear of
        attack.
\end{itemize}

\begin{DataFact}
National Highways are only \textbf{2\%} of India's road length, yet they carry about
\textbf{40\%} of the traffic and account for \textbf{30\%} of total fatalities.
\end{DataFact}

\begin{ExampleBox}
On the right to safe travel, the Court laid out three directions: (1) immediate
prohibition on the construction or operation of any new \emph{dhaba}, eatery or
commercial structure on the highway; (2) no heavy or commercial vehicle may park or stop
on a national highway carriageway or its paved shoulder; and (3) permission of a
\emph{district highway safety task force} is required. The reasoning is that roadside
structures and parked trucks block the right of way, and a driver travelling at 140--150
km/h cannot stop within 100 metres.
\end{ExampleBox}

\begin{QuickRevision}
The past year's Article 21 expansions are best remembered as a checklist: digital access
(Amar Jain), accessibility (Rajeev Raturi, 2024), trauma care (Save Life Foundation,
2026), timely pronouncement of reserved judgments (Peela Pahan), right to walk on a
footpath (Maniyar Iliyasi, Articles 19 and 21), safe travel on highways, and free
movement without fear of attack (the stray-dog/ABC-rules case).
\end{QuickRevision}

\section{Online Takedowns and Blocking versus Fundamental Rights}
\label{sec:online-blocking}

The second major block concerns the tension between online content takedowns / blocking
and fundamental rights under Article 19. The government passes blocking orders against
platforms such as Telegram and Facebook; the question is how far such orders are
justified against the background of internal security and public order.

\subsection{Statutory provisions}

The Information Technology Act, 2000 contains two central provisions:

\begin{itemize}
  \item \textbf{Section 69A} — blocking of access to information; it carries procedural
        safeguards and is therefore the more \emph{legally binding} provision.
  \item \textbf{Section 79(3)(b)} — the \emph{safe-harbour} clause for social media
        intermediaries; a conditional exemption from liability.
\end{itemize}

\begin{tabularx}{\textwidth}{YYY}
\toprule
\thead{Attribute} & \thead{Section 69A} & \thead{Section 79(3)(b)} \\
\midrule
Objective & Blocking of access to information & Safe-harbour protection to the
intermediary (conditional exemption from liability) \\
\addlinespace[3pt]
Grounds & Clearly stated: sovereignty and integrity of India, defence, security of the
state, friendly relations with foreign states, public order (similar to Article 19(2)) &
No separate grounds; per Shreya Singhal, must relate to the same grounds as 69A \\
\addlinespace[3pt]
Safeguards & Stringent — central government only, through the Ministry of Electronics and
IT & Lesser — other ministries plus states, through the Sahyog portal \\
\addlinespace[3pt]
Legal character & More legally binding & Due-diligence / safe-harbour provision \\
\bottomrule
\end{tabularx}

\subsection{The blocking procedure under Section 69A}

\begin{GovtPolicy}
Blocking under Section 69A operates through a \textbf{blocking committee}, chaired by a
\emph{designated officer} from the \textbf{Ministry of Electronics and IT}, with
representatives from the Ministry of Law and Justice, the Ministry of Home, the Ministry
of Information and Broadcasting, and the Computer Emergency Response Team (CERT-In). The
order must be approved by the \emph{Secretary}, Ministry of Electronics and IT. In an
\emph{emergency}, the designated officer may issue the order directly, provided it is
approved by the blocking committee \emph{within 48 hours}. The rules are the
\textbf{IT (Blocking) Rules, 2009} (distinct from the IT Rules, 2021, which deal with
significant social media intermediaries).
\end{GovtPolicy}

The \emph{grounds} of blocking are: sovereignty and integrity of India, defence, security
of the state, friendly relations with foreign states, and public order. What can be
blocked is \emph{access to any information generated, received, or hosted by any computer
resource}, adhered to by any agency of government or intermediary, and only through a
court order or an appropriate government notification.

\begin{CriticalPoint}
The designated officer (chairman of the blocking committee) is effectively ``just a
stamp.'' Since the 2009 rules were implemented, the blocking committee has \emph{never
once} disapproved a blocking order passed by the designated officer — which means there
is no parliamentary, judicial, or other safeguard that has ever reversed one of these
orders.
\end{CriticalPoint}

\subsection{The Telegram blocking order}

\begin{ExampleBox}
In the recent \textbf{Telegram blocking order}, the \emph{entire intermediary platform} of
Telegram was interpreted as ``information'' under Section 69A. The government took
recourse to Section 69A (rather than Section 79(3)(b)) because it is more legally
binding, gives the state a stronger hand in court, and was aimed at preventing paper
leaks. The Delhi High Court, on a petition by Telegram, held the blocking to be a
\emph{proportionate} exercise — proportionate to the objective of smooth conduct of the
examination and prevention of paper leakage. The law allows blocking at several levels: a
post, a user, a channel, or the whole platform; the criticism is that the state often
reaches for the platform-level step (``over-breadth'') when a narrower step would do.
\end{ExampleBox}

\textbf{Reasons for takedown.} The reasons the government invokes for takedown/blocking
orders are: (1) upholding national security and public order; (2) regulating cybercrime,
including malicious trading and dark-web URLs; (3) regulating VPNs; (4) regulating online
gaming and betting; and (5) prevention of paper leaks, piracy, and upholding the
intellectual property rights of creators, authors, producers and educators. To these one
adds the prevention of child pornography, for which the Court coined the term
\textbf{CSEM} — \emph{child sexual exploitative and abuse material}.

\begin{PYQLink}
The three questions the lecture framed for answer practice are: (1) ``Online digital free
speech is a vital expression of freedom of speech and expression but is subject to
regulation — discuss the need and issues concerning regulation under the IT Act, 2000.''
(2) ``Fundamental rights are not absolute — discuss in the light of recent
takedown/blocking orders of the government under the IT Act.'' (3) ``The IT Act, 2000 has
outlived its purpose — do you agree? Discuss with relevant case laws.''
\end{PYQLink}

\subsection{Issues concerning blocking and takedown orders}

\begin{itemize}
  \item \textbf{Restriction or violation of fundamental rights} — blocking affects the
        right to express, the right to receive information, and the right to criticise.
  \item \textbf{Over-breadth and the chilling effect} — account-level or platform-level
        blocking has a chilling effect on journalists, activists, political commentators,
        satirists, civil society organisations and media platforms.
  \item \textbf{Parallel exercise of Section 79(3)(b)} — the government continues to use
        Section 79 despite Section 69A being the statutory provision containing procedural
        safeguards.
  \item \textbf{Increasing censorship / decentralisation} — the government plans to
        decentralise blocking powers under Section 69A to several ministries, creating a
        censorship regime.
  \item \textbf{Dilution of Shreya Singhal (2015)} — unilateral decisions by the
        designated officer of the blocking committee, and the use of \textbf{Rule 16} of
        the Blocking Rules 2009 to keep blocking proceedings confidential (so the blocking
        order and its reasons are not provided to the affected parties), dilute the
        safeguards the Court relied on.
  \item \textbf{Lack of independent oversight} — there is no parliamentary oversight;
        blocking orders are reviewed only by the blocking committee, which is an executive
        body and has never overturned an order.
  \item \textbf{Limited legal remedy} — the mandatory hearing requirement under the
        blocking rules is bypassed by invoking emergency powers.
  \item \textbf{Collateral blocking} — banning a website or an intermediary also blocks
        \emph{lawful} content. Example: in July 2025, the blocking of accounts on X
        swept up the \emph{Reuters} account.
\end{itemize}

\begin{ExampleBox}
\textbf{Anuradha Bhasin (2020)} — in the context of the internet shutdown in Jammu and
Kashmir after the abrogation of Article 370, the Court held that freedom of speech and
expression, and the freedom to practise any profession \emph{over the internet}, are
protected by Articles 19(1)(a) and 19(1)(g).
\end{ExampleBox}

\subsection{Way forward}

The lecture's suggested reforms are:

\begin{enumerate}
  \item Account-level, domain-level, or URL-level blocking instead of platform-level
        blocking.
  \item An independent oversight mechanism led by the judiciary or Parliament.
  \item Exercise of blocking under Section 69A instead of Section 79.
  \item Publication of transparency reports by platforms — data on takedown requests,
        categories of content, and compliance rates.
  \item Time-bound orders with sunset clauses, to prevent permanent blocking.
  \item Regulation that also addresses AI-generated content (deepfakes and synthetically
        generated information), misinformation and fake news, and communal hatred.
\end{enumerate}

\begin{QuickRevision}
Section 69A is the blocking provision (with procedural safeguards and a blocking
committee chaired by a designated officer of MeitY, 48-hour emergency approval); Section
79(3)(b) is the safe-harbour clause. Shreya Singhal (2015) struck down 66A but upheld 69A
and 79 for their safeguards and judicial review; the Telegram order read the whole
platform as ``information'' and was upheld as proportionate by the Delhi High Court. The
criticisms are over-breadth/chilling effect, parallel use of 79, decentralisation,
dilution of Shreya Singhal, no independent oversight, limited remedy, and collateral
blocking (Reuters, July 2025). The remedy lies in narrower (URL-level) blocking,
independent oversight, 69A over 79, transparency reports, and sunset clauses.
\end{QuickRevision}

\section{Right to Bail}
\label{sec:right-to-bail}

The last block returns to a historically settled principle — \emph{bail is the rule, jail
is the exception} — now under stress from the Court's recent interpretation of stringent
laws.

\subsection{Constitutional basis}

\begin{GovtPolicy}
The constitutional basis of bail rests on four provisions:
\begin{itemize}
  \item \textbf{Article 21} — personal liberty; bail cannot become punishment before
        conviction.
  \item \textbf{Article 22} — safeguards against arbitrary arrest (for example, grounds
        of arrest must be communicated); these safeguards are essential to secure bail.
  \item \textbf{Article 14} — ensures non-arbitrariness in bail decisions.
  \item \textbf{Article 39A} — free legal aid.
\end{itemize}
\end{GovtPolicy}

\subsection{Statutory provisions}

The statutory framework is in the \textbf{Bharatiya Nyaya Sahita, sections 478--483},
which provide for bail and the release of undertrials after detention, bail in
non-bailable offences, anticipatory bail, and the special powers of the High Court and
Sessions Court.

\subsection{Important case laws}

\begin{itemize}
  \item \textbf{State of Rajasthan v. Balchand (1977)} — the classic principle that
        \emph{bail, not jail, is the rule}.
  \item \textbf{Hussainara Khatoon} — speedy trial, legal aid and undertrial liberty are
        integral to Article 21.
  \item \textbf{K.A. Najeeb (2021)} — statutory restrictions on bail under UAPA cannot
        justify \emph{indefinite incarceration} where the trial is unlikely to conclude
        soon.
\end{itemize}

\begin{CriticalPoint}
The recent interpretation (seen in the Andhrabi judgment, with Justice Bhuyan of Assam and
Justice B. Nagarathna) holds that even under stringent anti-terror laws such as
\textbf{UAPA} or \textbf{PMLA 2002}, prolonged incarceration — keeping a person behind
bars for five, six or seven years — becomes a ground for bail, because Articles 20 and 21
override the statutory framework of these anti-bail laws and restore the default that bail
is the rule and jail the exception. The Court noted that the courts themselves are at
divergence on this point.
\end{CriticalPoint}

\begin{QuickRevision}
Bail rests on Articles 21, 22, 14 and 39A, and is codified in BNS sections 478--483. The
case law is Balchand (1977, ``bail not jail''), Hussainara Khatoon (speedy trial and
undertrial liberty in Article 21), and K.A. Najeeb (2021, UAPA cannot justify indefinite
incarceration). The recent Andhrabi line holds that prolonged detention under stringent
laws itself becomes a ground for bail, because Articles 20 and 21 override the statutory
framework.
\end{QuickRevision}

\section*{Answer-Writing Pointers}
\addcontentsline{toc}{section}{Answer-Writing Pointers}

The lecture walked through one full mains question: \emph{``The Constitution is a living
document ensuring social justice. Discuss the statement in the light of recent Supreme
Court case laws.''}

\begin{PYQLink}
The three keywords of this question are \textbf{living document}, \textbf{social justice},
and \textbf{recent Supreme Court case laws}; a strong answer must ring all three bells.
\end{PYQLink}

The model opening line runs: the Constitution is a living document because the amendment
procedure under \textbf{Article 368} ensures flexibility along with rigidity, while the
provision of \textbf{judicial review} has led to an expansive framework of fundamental
rights. The body then moves in a fixed sequence:

\begin{enumerate}
  \item \textbf{Provisions of the Constitution ensuring social justice} — the Preamble,
        Fundamental Rights, DPSPs, and Articles 330/332 (reservation for women and
        Scheduled Castes/Tribes), along with recent Constitutional Amendment Acts: the EWS
        amendment (Articles 15(6) and 16(6)) and the 97th Amendment (2011, cooperatives).
  \item \textbf{The connecting link} — ``it is the judiciary which has played an important
        role in upholding constitutional morality.''
  \item \textbf{The judiciary's role, illustrated case-wise} — women (menstrual health,
        maternity/adoptive mothers), elderly, children, acid-attack victims, and the SC/ST
        sub-categorisation line (\textbf{Devinder Singh}), which held that homogenisation
        of Scheduled Castes is no longer a reality and permitted sub-categorisation for
        reservation.
  \item \textbf{Conclusion} — India's Constitution is based on a living ethos,
        \emph{constitutional morality}, or \emph{transformative constitutionalism}.
\end{enumerate}

\begin{CriticalPoint}
The answer-writing guidance is to absorb content first and then apply framing; the exam
rewards \emph{incremental changes} and the ability to put a specific case name to a
principle (e.g.\ Aligarh Muslim University) rather than restating common content.
\end{CriticalPoint}

\section*{Chapter Glossary}
\addcontentsline{toc}{section}{Chapter Glossary}

\begin{tabularx}{\textwidth}{YY}
\toprule
\thead{Term} & \thead{Meaning} \\
\midrule
Article 21 & Life and personal liberty; the main site of the Court's expansive
interpretation \\
Article 19(1)(a) & Freedom of speech and expression \\
Article 22 & Safeguards against arbitrary arrest \\
Articles 23--24 & Right against exploitation \\
Article 39A & Free legal aid (DPSP) \\
Article 39(e) & Protection of children from abuse (DPSP) \\
DPSP & Directive Principles of State Policy \\
Bail, not jail & Principle that bail is the rule and jail the exception \\
Anticipatory bail & Pre-arrest bail (BNS 478--483) \\
POCSO & Protection of Children from Sexual Offences Act, 2012 \\
MTP Act & Medical Termination of Pregnancy Act, 1971 \\
CARA & Central Adoption Resource Authority \\
Section 69A & Blocking of access to information (IT Act) \\
Section 79(3)(b) & Safe-harbour clause for intermediaries (IT Act) \\
Safe harbour & Conditional exemption of intermediaries from liability \\
Shreya Singhal & 2015 case striking down 66A, upholding 69A/79 safeguards \\
Chilling effect & Deterrence of lawful speech by over-broad blocking \\
Over-breadth & Reaching for platform-level blocking when a narrower step suffices \\
CSEM & Child sexual exploitative and abuse material \\
UAPA & Unlawful Activities (Prevention) Act \\
PMLA & Prevention of Money Laundering Act, 2002 \\
Judicial review & Power of courts to test legislative/executive action \\
Constitutional morality & The ethos/values underlying the Constitution \\
Transformative constitutionalism & Reading the Constitution to progressively advance rights \\
\bottomrule
\end{tabularx}
